\documentclass[10pt,a4paper]{amsart}
\usepackage[margin=19mm]{geometry}
\usepackage{amsmath,amssymb,mathtools}
\usepackage[scr=rsfs,cal=euler]{mathalfa}
\usepackage{xfrac}
\usepackage[unicode,hidelinks,bookmarks=false]{hyperref}
\newcommand{\ie}{i.e.\ }
\newcommand{\cf}{cf.\ }
\newcommand{\resp}{respectively\ }
\newcommand{\Subsection}{\subsection}
\providecommand{\nobreakdash}{\nobreak\hbox{-}}
\setlength{\emergencystretch}{2em}
\multlinegap0pt
\usepackage{amssymb}
\usepackage{mathtools}
\usepackage{bm}
\newtheorem{proposition}{Proposition}
\title{Exact Finite Integral Geometry: Directional Kernels and Spherical Designs}
\author{Congpei An}
\date{Source: arXiv:2609.05869v1 (CC BY 4.0)}
\begin{document}
\maketitle

\noindent\emph{Source and licence.} Exact Finite Integral Geometry: Directional Kernels and Spherical Designs. Version arXiv:2609.05869v1 (CC BY 4.0), distributed by arXiv under the Creative Commons Attribution 4.0 International licence, http://creativecommons.org/licenses/by/4.0/, as recorded in the arXiv licence metadata for this exact version and shown on the abstract page https://arxiv.org/abs/2609.05869v1. Full licence notice: \texttt{licenses/arxiv-2609.05869-CC-BY-4.0.txt}.\par\smallskip

\noindent\textit{Editorial scope.} Counted unit: the article's proposition prop:power-multipliers (source0.tex lines 458--468). The article's complete original proof of it is delivered with it (source0.tex lines 470--509, one \texttt{proof} environment of 40 lines). No mathematics is written, repaired or re-derived: the statement and the proof are the article's own lines, in the article's own notation, and no retained line is substituted or re-flowed.  Retained uncounted as the source-local context the proof needs: lines 453--456.  Nothing was omitted from any retained span, so the package carries no omission record.  No retained line contains a citation, so the wrapper carries no bibliography and no bibliography entry.  No retained line contains a figure environment, an image-inclusion command or drawing code, so no image is delivered and the package declares no asset dependency.  Editorial formatting only, in addition: an AMS \texttt{amsart} preamble that restates the article's own declarations and macros used by the retained lines, namely the environments \texttt{proposition} on separate counters, together with the standard packages those lines need.  The retained environments print in this document as Proposition 1 in order of appearance, whereas the article prints the same items with the numbering of its own full document.  The complete native source of the article as distributed in the arXiv e-print for this exact version is bundled unchanged as \texttt{sources/209469/source0.tex}.
\begin{equation}\label{eq:cdp}
 c_{d,p}:=\lambda_0(\psi_p)
 =\frac{\Gamma(d/2)\Gamma((p+1)/2)}{\sqrt\pi\,\Gamma((d+p)/2)}.
\end{equation}

\begin{proposition}[Funk--Hecke multipliers of $|s|^p$]\label{prop:power-multipliers}
Let $p>-1$.  For every $n\ge0$,
\begin{equation}\label{eq:power-multipliers}
 \lambda_{2n}(\psi_p)
 =c_{d,p}
 \frac{(-1)^n(-p/2)_n}{((d+p)/2)_n},
 \qquad
 \lambda_{2n+1}(\psi_p)=0.
\end{equation}
Here $(a)_n=\Gamma(a+n)/\Gamma(a)$ is the Pochhammer symbol.
\end{proposition}

\begin{proof}
Odd multipliers vanish because $\psi_p$ is even.  We give the Gegenbauer calculation for $d\ge3$ and note that the same formula in $d=2$ follows by the corresponding Fourier cosine integral.

Put $\alpha=(d-2)/2$.  With normalized spherical measure, the Funk--Hecke formula gives
\begin{equation}\label{eq:FH-gegen}
 \lambda_{2n}(\psi_p)
 =\frac{1}{B(1/2,\alpha+1/2)\,C_{2n}^{\alpha}(1)}
 \int_{-1}^1|s|^pC_{2n}^{\alpha}(s)(1-s^2)^{\alpha-1/2}\,ds.
\end{equation}
The quadratic change of variables $y=s^2$ and the identity
\begin{equation}\label{eq:gegen-jacobi}
 C_{2n}^{\alpha}(\sqrt y)
 =\frac{(\alpha)_n}{(1/2)_n}
 P_n^{(\alpha-1/2,-1/2)}(2y-1)
\end{equation}
reduce the integral to a beta--hypergeometric integral.  Using
\[
 P_n^{(\alpha-1/2,-1/2)}(2y-1)
 =(-1)^n\frac{(1/2)_n}{n!}
 {}_2F_1\!\left(\begin{matrix}-n,n+\alpha\\[1mm]1/2\end{matrix};y\right)
\]
and integrating the terminating series termwise yields a balanced ${}_3F_2(1)$.  Saalsch\"utz's summation gives
\begin{align}\label{eq:gegen-integral}
 &\int_{-1}^1|s|^pC_{2n}^{\alpha}(s)(1-s^2)^{\alpha-1/2}\,ds\\
 &\quad=
 B\!\left(\frac{p+1}{2},\alpha+\frac12\right)
 \frac{(-1)^n(\alpha)_n(\alpha+1/2)_n(-p/2)_n}
 {n!(1/2)_n(\alpha+p/2+1)_n}.
\end{align}
Finally,
\[
 C_{2n}^{\alpha}(1)
 =\frac{(\alpha)_n(\alpha+1/2)_n}{n!(1/2)_n},
\]
so substitution into \eqref{eq:FH-gegen}, together with
\[
 \alpha+\frac p2+1=\frac{d+p}{2},
\]
gives \eqref{eq:power-multipliers}.  The quotient of beta functions for $n=0$ is exactly \eqref{eq:cdp}.
\end{proof}

\end{document}
